The metrics of, \texttt{Total Analytical Trust Pre Experiment} and \texttt{Total Analytical Trust Post Experiment} are computed as row-wise sums of 10 analytical items at pre and post (Deceptive, Dishonest/Honest, Suspicious, Wary/Weary, Harm, Confident, Security, Trustworthy, Reliable, Trust).  
For analytical scoring, positively worded prompts are mapped from $-2$ (strongly disagree) to $+2$ (strongly agree), negatively worded prompts are reverse-coded, and higher totals therefore indicate greater analytical trust in AI. Each analytical item is coded in the data shared on Github and OSF with the analytical aligned original question and the scoring of the likert mapping. 

The metrics of \texttt{Total Emotional Trust Pre Experiment} and \texttt{Total Emotional Trust Post Experiment} are computed as row-wise sums of nine bipolar emotional prompts at pre and post: empathetic/apathetic, sensitive/insensitive, personal/impersonal, caring/ignoring, altruistic/self-serving, cordial/rude, responsive/indifferent, open-minded/judgmental, and patient/impatient. Each emotional item is coded in the data shared on Github and OSF with the trust-aligned adjective scored as $+1$ and the opposite adjective scored as $-1$. 

In all analysis the same scoring is used and scores are scaled as $A_{pre}=\frac{\text{Total Analytical Trust Pre Experiment}}{10}$, $A_{post}=\frac{\text{Total Analytical Trust Post Experiment}}{10}$, $E_{pre}=\frac{\text{Total Emotional Trust Pre Experiment}}{9}$, and $E_{post}=\frac{\text{Total Emotional Trust Post Experiment}}{9}$. Pre--post trust change is then defined as $\Delta A = A_{post}-A_{pre}$ and $\Delta E = E_{post}-E_{pre}$, where positive values indicate increased trust and negative values indicate decreased trust after the intervention. We also report overall trust as the simple mean of normalized analytical and emotional trust, with $O_{pre}=\frac{A_{pre}+E_{pre}}{2}$ and $O_{post}=\frac{A_{post}+E_{post}}{2}$. The pre--post overall trust change is then defined as $\Delta O = O_{post}-O_{pre}$.

\subsubsection{LLM Stimulus and Visualization Pipeline}
The experimental demonstrations were built around a masked language model based on BERT, specifically a version trained on Wikipedia and books that is described more fully in the codebase platform materials. Participant-facing examples use a WordPiece-style tokenizer with a vocabulary of roughly 30{,}000 subword units; input strings are lower-cased, segmented into subword tokens, and wrapped with the standard special markers $[CLS]$ and $[SEP]$ before inference. For the interactive fill-in-the-blank tasks, the tokenized sequence is passed to the completion pipeline, which returns the highest-probability candidate continuations and associated probabilities for the masked position; these predictions are then displayed as ranked token completions. The visual interface itself is rendered client-side with JavaScript and D3, combining HTML controls and SVG plots to show token-level completions, comparative sentence outputs, and temporal association graphics such as the gender-over-time panels. In this way, the study materials expose not only the model's textual predictions but also the relative structure of its learned associations across alternative prompts and demographic cues. All code for generating tokens and predictions of missing words in sentences is fully detailed on the online code repository. 

